% !TEX root = ../PRACA_MAIN.tex
\section*{Zakończenie} \addcontentsline{toc}{section}{Zakończenie}

Głównym problemem badawczym podjętym w niniejszej pracy było zjawisko przeciążenia poznawczego na platformach VOD oraz konieczność doboru optymalnych metod zautomatyzowanej personalizacji treści. Celem pracy było kompleksowe porównanie skuteczności różnorodnych algorytmów rekomendacyjnych, wywodzących się z trzech odmiennych paradygmatów: klasycznego filtrowania kolaboracyjnego, architektur uczenia głębokiego oraz nowoczesnych grafowych sieci neuronowych. Główna hipoteza badawcza zakładała, że zaawansowane modele uwzględniające nieliniowość i grafową topologię powiązań (takie jak NeuMF i LightGCN) wykażą znaczącą przewagę predykcyjną nad tradycyjnymi metodami opartymi na faktoryzacji macierzy i sąsiedztwie (ItemKNN, BPR), jednak odbędzie się to kosztem wymiernie wyższego narzutu obliczeniowego.

Przeprowadzone eksperymenty numeryczne na zbiorze danych MovieLens w pełni potwierdziły słuszność przyjętych założeń. Wyniki analiz jednoznacznie wykazały, że ewolucja architektur w kierunku bezpośredniego modelowania struktury grafu dwudzielnego przynosi wymierne korzyści w jakości rekomendacji. Model LightGCN osiągnął najwyższe wyniki w zestawieniu testowym pod kątem kluczowych metryk rangowych (w tym NDCG oraz MRR), udowadniając swoją wysoką zdolność do pozycjonowania najbardziej relewantnych obiektów na samym szczycie listy rekomendacyjnej. Z kolei hybrydowa architektura głęboka NeuMF okazała się najskuteczniejsza w ogólnym odnajdywaniu pożądanych pozycji, osiągając najwyższą wartość wskaźnika trafień. Jednocześnie zidentyfikowano i udokumentowano wyraźny kompromis pomiędzy precyzją układania rankingu a wymaganym czasem uczenia algorytmów. Nowoczesne sieci grafowe wymagały rzędy wielkości dłuższej optymalizacji niż klasyczne metody bazowe.

Biorąc pod uwagę uzyskane wyniki, można stwierdzić, że teza badawcza postawiona we wstępie pracy została udowodniona. Zrealizowano również główny cel badawczy, którym było praktyczne przetestowanie i porównanie wybranych algorytmów pod kątem ich trafności oraz obciążenia sprzętowego w kontekście platform streamingowych.

Wnioski płynące z przeprowadzonych badań posiadają istotne znaczenie praktyczne dla współczesnych przedsiębiorstw działających na rynku. Z biznesowego punktu widzenia wdrożenie najbardziej zaawansowanych algorytmów grafowych bez rzetelnej analizy opłacalności wiązałoby się ze znacznymi kosztami infrastrukturalnymi, które mogą być nieopłacalne w środowiskach wymagających błyskawicznego odświeżania bazy modeli. Przeprowadzona dyskusja wykazała, że w realiach ograniczonych zasobów sprzętowych starsza klasyczna faktoryzacja macierzy (BPR) lub głęboki model NeuMF stanowią wysoce pożądany balans pomiędzy satysfakcją końcowego odbiorcy a kosztami operacyjnymi utrzymania systemu.

Warto jednocześnie podkreślić pewne ograniczenia przeprowadzonych eksperymentów. Z uwagi na wysoką złożoność obliczeniową algorytmów opartych na grafowych sieciach neuronowych oraz ograniczone zasoby sprzętowe, analizę zrealizowano na podzbiorze MovieLens 100K. Wykorzystanie pełnoskalowych, komercyjnych zbiorów danych o znacznie większej rzadkości interakcji mogłoby dodatkowo uwydatnić różnice w precyzji pomiędzy architekturami, a jednocześnie wymusić zastosowanie bardziej wyrafinowanych strategii optymalizacyjnych niż klasyczne przeszukiwanie siatki. Kwestia ta stanowi naturalny kierunek dla przyszłych prac badawczych w tym obszarze.